\documentclass[11pt]{article}
\usepackage[margin=1in]{geometry}
\usepackage{booktabs}
\usepackage{longtable}
\usepackage{hyperref}
\usepackage{xcolor}
\usepackage{enumitem}
\title{Independent Replication Report\\
\large ``The Minimal Gene Complement of \textit{Mycoplasma genitalium}''\\
Fraser et al., \textit{Science} \textbf{270}, 397--403 (1995)}
\author{Ollie (OpenClaw AI) --- BVBRC-100 Replication Wave, target BVBRC-103}
\date{Report: 2026-07-04 \quad Backfilled to 8-artifact standard: 2026-07-05}
\begin{document}
\maketitle

\section{Paper summary}
Fraser et al.\ (1995) reported the second completely sequenced free-living bacterial genome,
\textit{Mycoplasma genitalium} G37, using whole-genome random-shotgun sequencing at TIGR
(following \textit{Haemophilus influenzae} Rd earlier the same year). The genome
(580{,}070~bp, single circular chromosome; GenBank L43967, now RefSeq NC\_000908.2) was at
publication the smallest known genome of any self-replicating organism, establishing
\textit{M.\ genitalium} as the paradigmatic minimal cellular life-form. The paper reported
$\sim$470 protein-coding ORFs, one rRNA operon (16S--23S--5S), 36 tRNAs covering all 20 amino
acids, and launched the ``minimal gene set'' program that culminated in JCVI-syn3.0 (2016).

\section{Claims tested}
\begin{longtable}{clcc}
\toprule
\# & Claim & Type & Tested? \\
\midrule
C1 & Genome length $=580{,}070$~bp, single circular chromosome & Quantitative & Yes \\
C2 & G+C content $\approx 32\%$ & Quantitative & Yes \\
C3 & $\sim$470 protein-coding genes & Quantitative & Yes (drift caveat) \\
C4 & One rRNA operon (16S+23S+5S) & Structural & Yes \\
C5 & 36 tRNAs & Quantitative & Yes \\
C6 & tRNAs cover all 20 amino acids & Structural & Yes \\
C7 & ``Smallest self-replicating genome'' (1995) & Historical & Yes \\
\bottomrule
\end{longtable}

\section{Method}
Real end-to-end reanalysis on free public data, local CPU only.
\textbf{Data:} NCBI E-utilities \texttt{efetch} of RefSeq NC\_000908.2 (GenBank flat file
780{,}009~B sha256 \texttt{50da1e36\ldots}; FASTA 588{,}426~B sha256 \texttt{cc21ace7\ldots}).
The GenBank record explicitly lists Fraser et al.\ 1995 as REFERENCE~2 for bases 1--580076,
confirming provenance. \textbf{Analysis:} Biopython 1.87 on Python 3.14 --- base composition,
feature-type counts, intact-vs-pseudogene CDS split, rRNA operon clustering (>5~kb gap = new
operon), tRNA amino-acid coverage, coding density, mean CDS length.
\textbf{Scoring:} LLM-as-judge via free Argo proxy (\texttt{argo:gpt-4o}, T=0.1) with per-claim
REPRODUCED/CLOSE/DIVERGENT judgments --- no regex scoring anywhere.

\section{Results vs paper}
\begin{longtable}{lrrl}
\toprule
Claim & Fraser 1995 & This work & Status \\
\midrule
C1 Genome length & 580{,}070~bp & 580{,}076~bp & REPRODUCED ($\Delta=+6$~bp, $10^{-5}$ rel.) \\
C2 G+C & $\sim$32\% & 31.69\% & REPRODUCED \\
C3 CDS & $\sim$470 & 504 intact + 20 pseudo & CLOSE (annotation drift) \\
C4 rRNA operons & 1 & 1 & EXACT \\
C5 tRNAs & 36 & 36 & EXACT \\
C6 20 aa & Yes & 20/20 & EXACT \\
C7 smallest (1995) & True & True & HISTORICALLY UPHELD \\
\bottomrule
\end{longtable}

\noindent Additional derived stats: coding density 93.04\%, mean CDS 356~aa,
A/T richness 68.31\%, N=0. LLM judge overall verdict: \textbf{REPLICATED}.

\section{Verdict}
\textbf{REPLICATED.} Every purely-sequence claim reproduces exactly or to sub-percent precision
with no parameter tuning. Data provenance is airtight (NC\_000908.2 cites Fraser 1995 as primary
reference). The only nontrivial delta --- 504 intact CDS vs $\sim$470 ORFs --- is expected
reannotation drift (PGAP gene-finders, overlapping-ORF resolution, comparative evidence), 7\%
relative, not touching the central near-minimal-genome finding.

\section{Genuine critique (evidence strength, shortcuts, unverified claims)}
\begin{itemize}[leftmargin=1.4em]
\item \textbf{This is not an \emph{ab initio} re-sequencing.} We re-analyzed the curated RefSeq
descendant of the Fraser sequence (6~bp corrected). We therefore verify \emph{claims about the
sequence}, not the \emph{shotgun-assembly process} that produced it in 1995. The assembly, gap
closure, and error-rate claims of the original paper are effectively untestable from public
artifacts and were \emph{not} checked.
\item \textbf{The C3 gene count comparison is apples-to-oranges by construction.} Fraser's
$\sim$470 came from 1995 ORF-finders; our 504 comes from 2020s PGAP. Calling this ``CLOSE'' is a
judgment, not a measurement --- we did \emph{not} rerun a period-appropriate gene-finder to
recover Fraser's actual 1995 call set, so we cannot state how many of the 470 survive vs how many
of the 504 are genuinely new. The ``annotation drift'' explanation is plausible but unproven here.
\item \textbf{Operon clustering used a heuristic threshold} (>5~kb gap = new operon). For a
single rRNA operon this is safe, but the method was not stress-tested against edge cases and
would need validation before use on genomes with multiple closely-spaced operons.
\item \textbf{LLM-judge scoring is a single-model, single-pass judgment} (argo:gpt-4o, T=0.1).
The standing policy prefers 3-judge scoring for robustness; this run used one judge. The verdict
is defensible because the underlying numbers are unambiguous, but the scoring layer is thinner
than policy ideal.
\item \textbf{Untested paper content:} minimal-gene-set derivation vs \textit{H.\ influenzae} Rd
(Fig.~2), energy-metabolism / cell-envelope / replication-repair completeness mapping, GC-skew,
and individual functional predictions. The replication is scoped to abstract-level quantitative
and structural claims only. Marking REPLICATED reflects those claims, not the full paper.
\end{itemize}

\section{Open questions}
See \texttt{report/open\_questions.json}. Five open problems surfaced: (Q1) essentiality
re-derivation vs modern Tn-seq; (Q2) exact 1995-vs-2026 gene-call reconciliation; (Q3)
minimal-genome comparison with JCVI-syn3.0; (Q4) pseudogene provenance; (Q5) codon/tRNA
sufficiency under the reduced set.

\end{document}
